\documentclass[a4paper]{article}
\usepackage{ctex}
\ctexset{
proofname = \heiti{证明}
}
\usepackage{amsmath, amssymb, amsthm}
\usepackage{moreenum}
\usepackage{mathtools}
\usepackage{url}
\usepackage{bm}
\usepackage{enumitem}
\usepackage{graphicx}
\usepackage{subcaption}
\usepackage{booktabs}
\usepackage[mathcal]{eucal}
\usepackage[thehwcnt = 3]{iidef}

\thecourseinstitute{自动化系}
\thecoursename{人工智能原理}
\theterm{2026年春季学期}
\hwname{作业}
\slname{\heiti{解}}

\begin{document}
\courseheader
\name{姓名：杜一郎\qquad \qquad \qquad 学号：2023011618}

\begin{enumerate}
\setlength{\itemsep}{3\parskip}

\item 4\(\times\)4 数独问题

已知数独盘面为
\[
\begin{matrix}
1 & 3 & \_ & 2\\
\_ & \_ & 3 & \_\\
\_ & 4 & \_ & 3\\
3 & \_ & 2 & \_
\end{matrix}
\]
要求每行、每列以及每个 \(2\times 2\) 子网格中均包含数字 1 至 4，且不重复。

\begin{enumerate}
\item 将该问题建模为约束满足问题，给出变量、值域和约束关系。
\item 求解该数独问题，给出关键步骤和填充后的完整数独网格。
\end{enumerate}

\begin{solution}
\textbf{(a) CSP 建模}

设第 \(i\) 行第 \(j\) 列格子对应变量为 \(X_{ij}\)，其中 \(i,j\in\{1,2,3,4\}\)。
因此变量集合为
\[
\mathcal X=\{X_{11},X_{12},\dots,X_{44}\}.
\]

每个变量的取值均来自集合 \(\{1,2,3,4\}\)。对于题目中已经给定的格子，其值域退化为单元素集合：
\[
X_{11}=1,\quad X_{12}=3,\quad X_{14}=2,\quad X_{23}=3,
\]
\[
X_{32}=4,\quad X_{34}=3,\quad X_{41}=3,\quad X_{43}=2.
\]
其余空格的初始值域均为
\[
D_{ij}=\{1,2,3,4\}.
\]

约束由数独规则给出，即每行、每列、每个 \(2\times 2\) 子网格中的四个变量满足 \texttt{Alldiff} 约束。

\textbf{行约束：}
\[
\mathrm{Alldiff}(X_{11},X_{12},X_{13},X_{14}),
\]
\[
\mathrm{Alldiff}(X_{21},X_{22},X_{23},X_{24}),
\]
\[
\mathrm{Alldiff}(X_{31},X_{32},X_{33},X_{34}),
\]
\[
\mathrm{Alldiff}(X_{41},X_{42},X_{43},X_{44}).
\]

\textbf{列约束：}
\[
\mathrm{Alldiff}(X_{11},X_{21},X_{31},X_{41}),
\]
\[
\mathrm{Alldiff}(X_{12},X_{22},X_{32},X_{42}),
\]
\[
\mathrm{Alldiff}(X_{13},X_{23},X_{33},X_{43}),
\]
\[
\mathrm{Alldiff}(X_{14},X_{24},X_{34},X_{44}).
\]

\textbf{子网格约束：}
\[
\mathrm{Alldiff}(X_{11},X_{12},X_{21},X_{22}),
\]
\[
\mathrm{Alldiff}(X_{13},X_{14},X_{23},X_{24}),
\]
\[
\mathrm{Alldiff}(X_{31},X_{32},X_{41},X_{42}),
\]
\[
\mathrm{Alldiff}(X_{33},X_{34},X_{43},X_{44}).
\]

\textbf{(b) 求解过程}

利用约束传播逐步缩减值域。

第 1 行为
\[
1,\ 3,\ \_,\ 2
\]
缺少的数字只有 4，因此
\[
X_{13}=4.
\]

左上 \(2\times 2\) 子网格为
\[
\begin{matrix}
1 & 3\\
\_ & \_
\end{matrix}
\]
缺少 \(\{2,4\}\)。而第 2 列目前为
\[
3,\ \_,\ 4,\ \_,
\]
缺少 \(\{1,2\}\)，所以
\[
X_{22}=2.
\]
进而左上子网格中另一个空格只能为 4，即
\[
X_{21}=4.
\]

第 2 行此时为
\[
4,\ 2,\ 3,\ \_,
\]
缺少的数字只有 1，因此
\[
X_{24}=1.
\]

第 1 列此时为
\[
1,\ 4,\ \_,\ 3,
\]
缺少的数字只有 2，因此
\[
X_{31}=2.
\]

第 3 行变为
\[
2,\ 4,\ \_,\ 3,
\]
缺少的数字只有 1，因此
\[
X_{33}=1.
\]

第 2 列为
\[
3,\ 2,\ 4,\ \_,
\]
缺少的数字只有 1，因此
\[
X_{42}=1.
\]

第 4 行变为
\[
3,\ 1,\ 2,\ \_,
\]
缺少的数字只有 4，因此
\[
X_{44}=4.
\]

故完整数独解为
\[
\begin{matrix}
1 & 3 & 4 & 2\\
4 & 2 & 3 & 1\\
2 & 4 & 1 & 3\\
3 & 1 & 2 & 4
\end{matrix}
\]
\end{solution}

\item 基站频率问题

6 个基站 \(A,B,C,D,E,F\) 呈环状分布。相邻两个基站交界处的观测强度等于这两个基站发射强度中的较大者。已知观测结果为
\[
\max(A,B)=3,\quad \max(B,C)=2,\quad \max(C,D)=2,
\]
\[
\max(D,E)=2,\quad \max(E,F)=1,\quad \max(F,A)=3.
\]

\begin{enumerate}
\item 将该问题建模为约束满足问题，给出变量、值域和约束关系。
\item 求解该问题，给出关键步骤和最终推测的基站信号强度等级。
\end{enumerate}

\begin{solution}
\textbf{(a) CSP 建模}

令变量集合为
\[
\mathcal X=\{A,B,C,D,E,F\},
\]
其中每个变量表示对应基站发射的信号强度等级。

每个变量的值域均为
\[
D_A=D_B=D_C=D_D=D_E=D_F=\{1,2,3\}.
\]

约束由题意“交界处观测强度等于相邻两个基站信号中的较大者”直接给出：
\[
\max(A,B)=3,
\]
\[
\max(B,C)=2,
\]
\[
\max(C,D)=2,
\]
\[
\max(D,E)=2,
\]
\[
\max(E,F)=1,
\]
\[
\max(F,A)=3.
\]

\textbf{(b) 求解过程}

由
\[
\max(E,F)=1
\]
可知两个变量都不能超过 1，因此必有
\[
E=1,\qquad F=1.
\]

再由
\[
\max(F,A)=3
\]
且 \(F=1\)，可得
\[
A=3.
\]

由
\[
\max(D,E)=2
\]
且 \(E=1\)，可得
\[
D=2.
\]

由
\[
\max(C,D)=2
\]
且 \(D=2\)，可知
\[
C\in\{1,2\}.
\]

由
\[
\max(B,C)=2
\]
可知 \(B\) 与 \(C\) 都不可能取 3，且二者至少有一个等于 2，因此
\[
B\in\{1,2\},\qquad C\in\{1,2\},\qquad \max(B,C)=2.
\]

最后检查约束
\[
\max(A,B)=3.
\]
由于 \(A=3\)，故无论 \(B=1\) 还是 \(B=2\)，该约束都自动满足。

因此最终可行解不是唯一的，而是共有以下三组：
\[
(A,B,C,D,E,F)=(3,2,1,2,1,1),
\]
\[
(A,B,C,D,E,F)=(3,1,2,2,1,1),
\]
\[
(A,B,C,D,E,F)=(3,2,2,2,1,1).
\]

故该题全部可行解如上。
\end{solution}

\item 生产排期问题

四个定制产品 \(A,B,C,D\) 需要共用一套生产设备。可用时段为
\[
1:\text{周一上},\qquad 2:\text{周一下},\qquad 3:\text{周二上},\qquad 4:\text{周二下}.
\]
约束条件如下：
\begin{itemize}
  \item 产品 \(A\) 和产品 \(B\) 不能安排在相邻时段；
  \item 产品 \(C\) 只能在下午生产；
  \item 产品 \(D\) 必须先于产品 \(A\) 完成生产；
  \item 如果产品 \(B\) 在周二生产，那么产品 \(C\) 必须在周一完成生产。
\end{itemize}

\begin{enumerate}
\item 将该问题建模为约束满足问题，给出变量、值域和约束关系。
\item 求解该问题，给出关键步骤和最终生产排期方案。
\end{enumerate}

\begin{solution}
\textbf{(a) CSP 建模}

令变量集合为
\[
\mathcal X=\{A,B,C,D\},
\]
其中每个变量的值表示该产品被安排的时段编号。

初始值域均为
\[
D_A=D_B=D_C=D_D=\{1,2,3,4\}.
\]

由“产品 \(C\) 只能在下午生产”可得一元约束
\[
C\in\{2,4\}.
\]
因此
\[
D_C=\{2,4\}.
\]

其余约束如下：

\begin{enumerate}[label=\arabic*)]
  \item 四个产品占用独立时段，即
  \[
  \mathrm{Alldiff}(A,B,C,D).
  \]

  \item 产品 \(A\) 和产品 \(B\) 时段不能相邻，即
  \[
  |A-B|\neq 1.
  \]

  \item 产品 \(D\) 必须先于产品 \(A\)，即
  \[
  D<A.
  \]

  \item 若 \(B\) 在周二生产，则 \(C\) 必须在周一完成生产。周二对应时段 \(3,4\)，周一对应时段 \(1,2\)，因此该约束可写成
  \[
  B\in\{3,4\}\Rightarrow C\in\{1,2\}.
  \]
  结合 \(C\in\{2,4\}\)，可进一步理解为：若 \(B\in\{3,4\}\)，则必有 \(C=2\)。
\end{enumerate}

\textbf{(b) 求解过程}

由 \(C\in\{2,4\}\)，分两种情况讨论。

\textbf{情况一：}\(C=4\)。

此时若 \(B\in\{3,4\}\)，根据条件“若 \(B\) 在周二生产，则 \(C\) 必须在周一完成”，会推出 \(C\in\{1,2\}\)，与 \(C=4\) 矛盾。因此
\[
B\notin\{3,4\},
\]
即
\[
B\in\{1,2\}.
\]

由于四个产品时段互不相同，而 \(C=4\)，所以 \(A,B,D\) 必须占用时段 \(\{1,2,3\}\)。

若 \(B=1\)，则由于 \(A\) 与 \(B\) 不能相邻，故 \(A\neq 2\)，又 \(A\neq 1\)，只能有
\[
A=3.
\]
于是剩余时段给 \(D\)，得
\[
D=2.
\]
并且满足 \(D<A\)。故得到一组可行解
\[
(A,B,C,D)=(3,1,4,2).
\]

若 \(B=2\)，则 \(A\) 不能取 1，也不能取 3，因为这两个时段都与 2 相邻，因此无解。

所以在 \(C=4\) 的情况下，唯一可行解为
\[
(A,B,C,D)=(3,1,4,2).
\]

\textbf{情况二：}\(C=2\)。

此时 \(A,B,D\) 需要从时段 \(\{1,3,4\}\) 中取值，且满足 \(D<A\)。

先讨论 \(A\) 的可能取值。

若 \(A=1\)，则不可能存在 \(D<A\)，故无解。

若 \(A=3\)，则只能有 \(D=1\)，剩余时段给 \(B\)，得 \(B=4\)。但此时 \(A=3\) 与 \(B=4\) 相邻，矛盾，因此无解。

若 \(A=4\)，则 \(D\) 可取 1 或 3：
\begin{itemize}
  \item 若 \(D=1\)，则 \(B=3\)。但 \(A=4\) 与 \(B=3\) 相邻，故无解；
  \item 若 \(D=3\)，则 \(B=1\)。此时 \(|A-B|=|4-1|=3\neq 1\)，且 \(D<A\) 成立，因此得到可行解
  \[
  (A,B,C,D)=(4,1,2,3).
  \]
\end{itemize}

综上，全部可行排期方案共有两组：
\[
(A,B,C,D)=(3,1,4,2),
\]
\[
(A,B,C,D)=(4,1,2,3).
\]

将其翻译成具体时间安排，可写为：
\begin{enumerate}[label=方案\arabic*：]
  \item \(B\) 在周一上，\(D\) 在周一下，\(A\) 在周二上，\(C\) 在周二下；
  \item \(B\) 在周一上，\(C\) 在周一下，\(D\) 在周二上，\(A\) 在周二下。
\end{enumerate}
\end{solution}

\end{enumerate}

\end{document}
